% preamble.tex — Companion paper: Finite Primitive Basis Theorem
% Journal-style formatting for the formal foundations paper

\documentclass[11pt]{article}

%% ——— Geometry ———
\usepackage[
  top=2.54cm,
  bottom=2.54cm,
  left=2.54cm,
  right=2.54cm
]{geometry}

%% ——— Spacing ———
\usepackage{setspace}
\onehalfspacing

%% ——— Mathematics ———
\usepackage{amsmath}
\usepackage{amssymb}
\usepackage{amsthm}

%% ——— Colors (must load before pgfplots/tikz to avoid option clash) ———
\usepackage[dvipsnames]{xcolor}

%% ——— Graphics and figures ———
\usepackage{graphicx}
\usepackage{tikz}
\usepackage{pgfplots}
\pgfplotsset{compat=1.18}

%% ——— Tables ———
\usepackage{booktabs}
\usepackage{multirow}
\usepackage{array}
\usepackage{longtable}

\definecolor{pwmBlue}{HTML}{2E5FA1}
\definecolor{pwmOrange}{HTML}{E8712B}
\definecolor{pwmGreen}{HTML}{3D9E3F}
\definecolor{pwmRed}{HTML}{C93C3C}

%% ——— Hyperlinks and cross-references ———
\usepackage[
  colorlinks=true,
  linkcolor=pwmBlue,
  citecolor=pwmBlue,
  urlcolor=pwmBlue
]{hyperref}
\usepackage[capitalise,noabbrev]{cleveref}

%% ——— Bibliography ———
\usepackage[numbers,sort&compress]{natbib}

%% ——— Theorem environments ———
\newtheorem{theorem}{Theorem}
\newtheorem{lemma}[theorem]{Lemma}
\newtheorem{proposition}[theorem]{Proposition}
\newtheorem{corollary}[theorem]{Corollary}

\theoremstyle{definition}
\newtheorem{definition}[theorem]{Definition}
\newtheorem{example}[theorem]{Example}
\newtheorem{remark}[theorem]{Remark}

%% ——— Custom commands ———
\newcommand{\pwm}{\textsc{pwm}}
\newcommand{\opg}{\textsc{OperatorGraph}}
\newcommand{\ctier}{\mathcal{C}_{\mathrm{Tier2}}}
\newcommand{\Blib}{\mathcal{B}}
\newcommand{\Hnom}{H_{\mathrm{nom}}}
\newcommand{\Htrue}{H_{\mathrm{true}}}
\newcommand{\Hdag}{H_{G}}
\newcommand{\etier}{e_{\mathrm{Tier2}}}
\newcommand{\Nmax}{N_{\max}}
\newcommand{\Dmax}{D_{\max}}
\newcommand{\compose}{\mathrm{compose}}
\newcommand{\ie}{\textit{i.e.}}
\newcommand{\eg}{\textit{e.g.}}
